\clearpage
\SourceLeaf{402}\pdfbookmark[1]{II. Tables for the Perturbations of Ceres}{ceres-tables-letters}
\begin{center}
{\large CORRESPONDENCE.\par}
\bigskip
[II. TABLES FOR THE PERTURBATIONS OF CERES.]\par
\bigskip
\textsc{Gauss} to \textsc{Olbers}, Braunschweig, 3 December 1802.
\end{center}
\noindent\hspace*{\parindent}.\quad.\quad.\quad.\quad.\quad I am now computing a table for the perturbations that the \rlap{\kern.02em\raisebox{1.02em}{\rule{.2em}{.25pt}}}Nov. issue of the M. C. r\rlap{\kern.35em\raisebox{1.02em}{\rule{.2em}{.25pt}}}eports for {\AstroSymbols ⚳}. For part of it I am using an arrangement that differs from the usual one. The tables for the perturbations in longitude and radius vector that depend only on \(\text{\AstroSymbols ⚳}-\text{\AstroSymbols ♃}\), and those for latitude, are already finished. I will write to you about it in more detail soon.\quad.\quad.\quad.\quad.\quad.

\begin{center}
\textsc{Gauss} to \textsc{Olbers}, Braunschweig, 21 December 1802.
\end{center}
\noindent\hspace*{\parindent}.\quad.\quad.\quad.\quad.\quad My now completed tables\refstepcounter{footnote}\hyperlink{ceresII-note}{[*)]}\begingroup
\makeatletter\long\def\@makefntext#1{\noindent[*) #1}\makeatother
\footnotetext{\hypertarget{ceresII-note}{}Vol. VI, pp. 235—243.]}
\endgroup{} for the perturbations of {\AstroSymbols ⚳} by {\AstroSymbols ♃} will go to \textsc{Zach} this week. If they are still too extensive for the M. C., I will get you a copy. I can describe the arrangement to you in a few words. For longitude, first, all the equations whose argument is a multiple of \(\text{\AstroSymbols ⚳}-\text{\AstroSymbols ♃}\) make up one table. The sum of all the remaining equations that depend on the first power of the eccentricity can be expressed by \(A\sin(B-\text{\AstroSymbols ⚳})\), so that \(A\) and \(B\) are functions only of \(\text{\AstroSymbols ⚳}-\text{\AstroSymbols ♃}\). A second table gives these quantities \(A\) and \(B\) for every value of \(\text{\AstroSymbols ⚳}-\text{\AstroSymbols ♃}\), degree by degree. In the same way there are 2 tables for the radius vector and one for the latitude, 5 tables in all, whereas one would otherwise need 40 tables for all the equations I have used, or still some 30 \SourceLeaf{403}if one omitted those that do not reach \(2''\). For me at least it is \textls[180]{far} easier to do these simple trigonometric calculations than to form so many arguments, consult so many tables and perform so many additions. My table could easily be made still easier to use by replacing \(A\) in the table with \(\log A\), for example. Alternatively, one could choose a constant quantity \(=2M\), no smaller than any value of \(A\), and set \(\frac{A}{M}=\cos C\). The sum of the equations that depend on the eccentricity would then be
\[
=M\sin(B-C-\text{\AstroSymbols ⚳})+M\sin(B+C-\text{\AstroSymbols ⚳}).
\]
For each value of \(\text{\AstroSymbols ⚳}-\text{\AstroSymbols ♃}\), one would then enter the values of \(B-C\) and \(B+C\) in the table and add a table giving \(M\sin\varphi\) for all values of \(\varphi\). — Incidentally, I have used, in place of the angles \(B\), the following values for
\[
\begin{array}{@{}l@{\quad}l@{}}
\text{longitude}\quad\text{.\ .\ .\ .}&B-2(\text{\AstroSymbols ⚳}-\text{\AstroSymbols ♃})\\[3pt]
\left.\begin{array}{@{}l@{}}\text{the R. V.}\\\text{latitude}\end{array}\right\}\ \text{.\ .\ .}&B-3(\text{\AstroSymbols ⚳}-\text{\AstroSymbols ♃})
\end{array}
\]
or given the sum of the equations the form
\[
\begin{array}{@{}l@{\quad}l@{}}
\text{for longitude}\ \text{.\ .\ .\ .\ .\ .\ .\ .\ .\ .\ .\ .\ .\ .\ .}&A\sin(B'+\text{\AstroSymbols ⚳}-2\text{\AstroSymbols ♃})\\[4pt]
\text{for the R. V. and lat.}\ \text{.\ .\ .\ .}&A\sin(B'+2\text{\AstroSymbols ⚳}-3\text{\AstroSymbols ♃})
\end{array}
\]
in my table.

This \(B'\) thus always remains within certain bounds as \(\text{\AstroSymbols ⚳}-\text{\AstroSymbols ♃}\) runs through all its values, whereas \(B\) would make one or two revolutions during one period of \(\text{\AstroSymbols ⚳}-\text{\AstroSymbols ♃}\).\quad.\quad.\quad.\quad.\quad.

\begin{center}
\textsc{Gauss} to \textsc{Olbers}, Braunschweig, 4 January 1803.
\end{center}
\noindent\hspace*{\parindent}.\quad.\quad.\quad.\quad.\quad I have already made many calculations with my provisional tables for the {\AstroSymbols ⚳} perturbations .\quad.\quad.\quad.\quad and found their exceptional \textls[180]{convenience} borne out. But I am still not quite satisfied with the VIII. set of elements printed in the Nov. issue of the M. C. For these I had used meridian observations alone. These were represented very well by that set, and much better than by the\,\raisebox{.65ex}{.} VII. set. But now that the perturbation calculations, so \SourceLeaf{404}laborious without tables, no longer pose an obstacle, I have also compared \textsc{Oriani}'s latest observations with those elements and, to my surprise, found the error appreciably larger than with the VII. elements. The right ascensions all agree fairly well, but the declinations are quite consistently about \(40''\) too small. I therefore put the error in longitude for 5 Aug. at \(+18''\), and the error in latitude at \(-37''\) on average. The error of the VII. elements is only about half as large. It seems to me that this confirms an opinion I expressed some time ago. The uncertainty resulting from the still insufficient span of the observations may for now impair predictions of the future motion more than neglecting the perturbations. I therefore intend to refine the elements further and fit them as closely as possible to all the available observations. I would in any event have to make nearly all the calculations this requires in this year, 1803, to make use of future observations. The work is therefore well spent, though it is hardly necessary for recovering the body in the eastern sky. For that purpose, I think, the ephemerides calculated by \textsc{Triesnecker} and \textsc{Bode} will be sufficiently accurate. — What chiefly still prevents great precision in the elements is the short span of \textsc{Piazzi}'s observations from 1801. For an accurate determination, when as here the mean motion cannot be derived from observations far apart in time, it is \textls[180]{essential} to have 4 observations \(A,B,C,D\), all sufficiently far apart, and if possible to make the interval \(B,C\) much larger than \(A,B\) and\,\raisebox{.8ex}{.} \(C,D\). For this one has no choice but to take as \(A\) and \(B\) the first and last of \textsc{Piazzi}'s observations, which are only 41 days apart. One can therefore expect a much greater degree of accuracy when one needs only one of \textsc{Piazzi}'s observations for \(A\) (which one of course makes more accurate than a single observation by comparing it with the neighboring ones), that is, in 1803,
\[\textstyle
\text{when }\left\{\begin{array}{c}A\\B\\C\\D\end{array}\right\}
\text{ from }\left\{\begin{array}{ll}
1801&\\
1802&\text{(or \textsc{Zach}'s obs. of 7 Dec. 1801)}\\
1802&\\
1803&
\end{array}\right.
\]
\SourceLeaf{405}can be taken. But one can only hope for really high accuracy after the {\AstroSymbols ☍} of 1804, when one can take \(A,B,C,D\) from as many different years. In 1806, when one can use 4 oppositions, we will hopefully know the orbit of {\AstroSymbols ⚳} almost as accurately as those of the other planets, with the possible exception of the mean motion. With {\AstroSymbols ⚴}, everything will take one year longer. — But if we find one of the two again in an older observation, a hope we have yet to give up, everything will of course proceed much faster.\quad.\quad.\quad.\quad.\quad.

\begin{center}
\textsc{Gauss} to \textsc{Olbers}, Braunschweig, 1 March 1803.
\end{center}
\noindent\hspace*{\parindent}.\quad.\quad.\quad.\quad.\quad I would have sent you my tables for the perturbations of {\AstroSymbols ⚳} long ago, had von \textsc{Zach} not written to me that everything would appear in the February issue of the M. C. As there will not have been enough room in that issue, you will receive them with my next letter. I have also computed with little effort the other tables needed for the motion of {\AstroSymbols ⚳}, to minutes only, really just for my own use in examining the Histoire Céleste that I received very recently. I have already looked through it and made a few quick trials. But I almost despair of finding {\AstroSymbols ⚳} in it, although I came across several days on which the difference in declination was only a few degrees. I will send you this table too before long. Both tables are only provisional, however. After this year's observations are complete, we will have far more accurate elements, and it will then be worth the effort to compute and tabulate the perturbations more accurately and completely.\quad.\quad.\quad.\quad.\quad.

\begin{center}
\textsc{Gauss} to \textsc{Olbers}, Braunschweig, 8 April 1803.
\end{center}
\noindent\hspace*{\parindent}.\quad.\quad.\quad.\quad.\quad I am sending you herewith the long-promised small tables for {\AstroSymbols ⚳} and {\AstroSymbols ⚴}. You will receive the one still outstanding for {\AstroSymbols ⚴} shortly.\quad.\quad.\quad.\quad.\quad.

Only this year will I be able to apply to the further improvement of the {\AstroSymbols ⚳} elements the method that can then be used for continued refinement. Until now it was too early for that. As this \SourceLeaf{406}method is in my opinion a significant part of the whole, I intend to take up the treatment of the whole once I have tested this part in \textls[180]{practice} and, as usually happens, perfected it in practice.\quad.\quad.\quad.\quad.\quad.

\begin{center}
\textsc{Gauss} to \textsc{Olbers}, Braunschweig, 25 January 1805.
\end{center}
\noindent\hspace*{\parindent}.\quad.\quad.\quad.\quad.\quad My new elements for Ceres .\quad.\quad.\quad.\quad.\quad are as follows.
\begin{center}
\begingroup\small\setlength{\tabcolsep}{5pt}\renewcommand{\arraystretch}{1.25}
\begin{tabular}{r r r r}
 & Epoch & Aphelion & Eccentricity\\
1801 & \(77^0\ 17'\ \ 0''{,}1\) & \(326^0\ 19'\ 59''\) & \(0{,}078\ 4929\)\\
1802 & \(155\ \ 27\ \ 34{,}2\) & \(22\ \ \ 0\) & \(4871\)\\
1803 & \(233\ \ 38\ \ \ 8{,}3\) & \(24\ \ \ 2\) & \(4814\)\\
1804 & \(312\ \ \ 1\ \ 33{,}5\) & \(26\ \ \ 3\) & \(4757\)\\
1805 & \(30\ \ 12\ \ \ 7{,}7\) & \(28\ \ \ 4\) & \(4700\)
\end{tabular}\endgroup
\end{center}
\begin{center}
\begingroup\footnotesize\setlength{\tabcolsep}{2pt}\renewcommand{\arraystretch}{1.25}
\begin{tabular}{l@{\quad}r}
Logarithm of the semimajor axis .\ .\ .\ .\ . & \(0{,}442\ 0004\)\\
Daily tropical motion .\ .\ .\ .\ .\ .\ .\ .\ .\ . & \(771''{,}0524\)\\
\(\AscendingNode\) 1801 .\ .\ .\ .\ .\ .\ .\ .\ .\ .\ .\ .\ .\ .\ .\ .\ .\ .\ .\ .\ .\ .\ .\ .\ .\ .\ .\ . & \(80^0\ 54'\ 46''\)\\
Inclination 1801 .\ .\ .\ .\ .\ .\ .\ .\ .\ .\ .\ .\ .\ .\ .\ .\ .\ .\ .\ .\ .\ .\ . & \(10^0\ 38'\ 13''\).
\end{tabular}\endgroup
\end{center}

These elements, which must still be used together with my old perturbation tables, are based on the 3 {\AstroSymbols ☍} and \textsc{Piazzi}'s observations from 1801. Yet the \textls[180]{latitudes} in the 4 years cannot be represented by the secular change found for \(\AscendingNode\) and inclination (I already noticed the former last year). To reconcile the observations with my perturbation tables I must assign
\begin{center}
\begin{tabular}{l l l}
to the node & \(0''{,}243\) & daily displacement\\
to the inclination & \(0''{,}025\) & daily decrease
\end{tabular}
\end{center}
, both much larger than the calculation gives. No doubt this is simply due to the incompleteness of the periodic equations for the radius vector and the latitude, and will become clear in a more elaborate calculation of the perturbations.
\par\medskip\begin{center}\rule{.22\linewidth}{.4pt}\end{center}
